\section{Введение}
Астроцитарная кальциевая сигнализация связывает биохимические процессы внутри клетки с её взаимодействиями с нейронами. Включение таких процессов в сетевую модель требует выбора между подробностью описания отдельного астроцита и вычислительной стоимостью всей системы. Разнообразие существующих кальциевых моделей и различия в их назначении подробно рассмотрены в обзоре~\cite{Manninen2018}. Возможность распределённого моделирования нейроастроцитарных сетей, включая сети до миллиона клеток, показана в работе~\cite{Jiang2025}. Следовательно, вопрос состоит не в принципиальной невозможности больших расчётов, а в стоимости выбранного уровня детализации и в том, какие свойства сохраняет более компактная модель.

Один из путей снижения этой стоимости --- феноменологическое описание: вместо воспроизведения всех биохимических потоков модель должна сохранять выбранные типы равновесий, механизмы возбуждения и характеристики кальциевых колебаний. В нейронном моделировании известным примером такого подхода служит модель Ижикевича~\cite{Izhikevich2003}. Для астроцитов полезно ставить аналогичную задачу построения компактного вычислительного элемента. Однако отсутствие единственной общепринятой модели не следует понимать как отсутствие редукций или аппаратных реализаций. В частности, предложены цифровые модели астроцитов и нейроглиальных взаимодействий для FPGA~\cite{DigitalAstrocyte2015}.
\todo{Дополнить обзор конкретными работами по феноменологической редукции и сравнением сохраняемых ими свойств. Утверждение об отсутствии общепринятого аналога модели Ижикевича пока является мотивацией, а не результатом систематического обзора.}

Непрерывные кусочно-аффинные (PWA) системы представляют интерес для этой задачи по двум причинам. Во-первых, изменение локального спектра между областями позволяет воспроизводить различные механизмы возникновения колебаний. Во-вторых, внутри области решение задаётся матричной экспонентой, что позволяет отделить распространение состояния от поиска следующего пересечения границы. Теория Hopf-подобных переходов и предельных циклов таких систем развита в работах~\cite{Simpson2022,SimpsonMeiss2007,LlibrePonceValls2015,ChenLiXieYue2017,ChenFengJiaMeng2026}. Аналитичность регионального потока при этом не означает гладкости полного поля на границах или явной формулы для времени каждого переключения.

В настоящей работе исследуется трёхзонный осциллятор с двумя порогами, ориентированный на кальциевые ответы Li--Rinzel. Сначала определяются свойства эталонной динамики, которые требуется сохранить. Затем рассматриваются локальные и большие циклы PWA-системы, вводится параметризация модели и проводится сравнение вариантов подбора. Наконец, точность кальциевых ответов сопоставляется с вычислительными затратами в NEST, где восстановленный кальций используется для вычисления медленного входящего тока SIC.

\section{Эталонная динамика Li--Rinzel}
\subsection{Уравнения и режимы модуляции}
В качестве эталона используется двухпеременная модель~\cite{LiRinzel1994,DePitta2009}
\begin{align}
\dot C&=\left[r_C m_\infty^3 n_\infty^3h^3+r_L\right]
       [C_0-(1+c_1)C]\notag\\
       &\quad-v_{ER}\frac{C^2}{C^2+K_{ER}^2},\\
\dot h&=a_2[Q_2-(Q_2+C)h],\\
m_\infty&=\frac{I}{I+d_1},\qquad n_\infty=\frac{C}{C+d_5},\\
Q_2&=d_2\frac{I+d_1}{I+d_3}.
\end{align}
Здесь $C$ и $I$ --- концентрации кальция и IP$_3$ в мкМ, $h$ --- безразмерная переменная доступности рецептора, время измеряется в секундах. Использованы $C_0=2$, $c_1=0.185$, $r_C=6$, $r_L=0.11$, $v_{ER}=0.9$, $a_2=0.2$, $d_1=0.13$, $d_2=1.049$, $d_3=0.9434$ и $d_5=0.08234$ в согласованных единицах. Рассматриваются $K_{ER}=0.1$, $0.075$ и $0.05$ мкМ, обозначаемые далее AM, AFM и FM. Эти обозначения характеризуют изменение амплитуды и периода в исследуемом диапазоне входов, а не строго непересекающиеся классы.

При постоянном $I$ нульклина рецепторной переменной имеет вид $h=Q_2/(Q_2+C)$; кальциевая нульклина находится из $\dot C=0$. Равновесия, экстремумы периодических орбит и их периоды представлены на рис.~\ref{fig:ru-lr}. Устойчивость равновесий определяется собственными числами якобиана, а устойчивость циклов --- нетривиальным множителем Флоке. Наличие устойчивого цикла устанавливается отдельно от неустойчивости равновесия.
\rufigure{figure1_li_rinzel}{Кальциевые колебания Li--Rinzel при постоянном IP$_3$: AM, AFM и FM сверху вниз. Слева показаны нульклины и устойчивые циклы, в центре --- равновесия и экстремумы циклов, справа --- периоды. Заполненные и пустые маркеры соответствуют устойчивым и неустойчивым циклам. Квадраты обозначают точки Хопфа, ромбы --- верхние складки циклов. Все кривые рассчитаны независимо; композиция следует Fig.~2 работы~\cite{DePitta2009}. Длиннопериодные FM-ветви ограничены точностью независимой проверки замыкания; крайняя точка не отождествляется с глобальной бифуркацией.}{fig:ru-lr}{}

\subsection{Спектральные требования к редуцированной модели}
Для якобиана $J$ введём след $T=\operatorname{tr}J$, определитель $D=\det J$ и дискриминант $\Delta=T^2-4D$. При $D>0$, $T<0$ равновесие асимптотически устойчиво: оно является фокусом при $\Delta<0$ и узлом при $\Delta\geq0$, с отдельным рассмотрением вырожденных случаев. Центральный колебательный участок эталона и верхнее устойчивое состояние мотивируют фокусную динамику соответствующих областей PWA-модели. Для низкоактивного состояния фиксируется устойчивость, но допускаются узел и фокус. Центр не включается в это условие, поскольку не является асимптотически устойчивым.

Такое требование описывает целевую организацию региональных полей. Оно не предполагает, что все три региональных равновесия одновременно принадлежат своим областям. Кроме того, в FM-эталоне имеются седловые участки, которые нельзя исключать при заявлении о полном совпадении структур. Поэтому речь идёт о феноменологическом воспроизведении выбранных режимов, а не о доказанной топологической эквивалентности Li--Rinzel и редуцированной системы.
\rufigure{figure2_lr_spectrum}{Спектральная карта равновесий Li--Rinzel в координатах след--определитель. Фон разделяет узлы, фокусы и седла. Границы $D=0$, $T=0$ и $D=T^2/4$ задают спектральные вырождения; для установления бифуркации требуются дополнительные условия. Параметризация по равновесному кальцию сохраняет складчатые ветви.}{fig:ru-spectrum}{}

\section{Трёхзонная модель и механизмы колебаний}
\subsection{Параметризация}
Подбираемые модели записываются в координатах $(z,q)$:
\begin{equation}\label{eq:pwa-ru}
\begin{gathered}
\dot z=a[q-N(z,I)],\quad \dot q=b[H(z,I)-q],\\ a,b>0,
\end{gathered}
\end{equation}
где
\begin{align}
N(z,I)&=\alpha_0+\alpha_1z+\alpha_I I\notag\\
 &\quad+\alpha_L(z-z_L)_++\alpha_R(z-z_R)_+,\\
H(z,I)&=\beta_0+\beta_1z+\beta_I I\notag\\
 &\quad+\beta_L(z-z_L)_++\beta_R(z-z_R)_+,
\end{align}
$(u)_+=\max(u,0)$ и $z_L<z_R$. Общая шарнирная запись обеспечивает непрерывность поля. Две прямые разделяют плоскость на три области с постоянными матрицами при фиксированных коэффициентах. В каждой области $N=n_jz+n_{0j}+\alpha_I I$ и $H=h_jz+h_{0j}+\beta_I I$, поэтому
\begin{equation}
\begin{gathered}
J_j=\begin{pmatrix}-a n_j&a\\ b h_j&-b\end{pmatrix},\\
T_j=-a n_j-b,\qquad D_j=ab(n_j-h_j).
\end{gathered}
\end{equation}
Налагаются $D_j>0$; для средней области $T_j>0$, $T_j^2<4D_j$, для правой $T_j<0$, $T_j^2<4D_j$, для левой $T_j<0$ без обязательного ограничения знака дискриминанта. Равновесие каждой аффинной подсистемы является действительным равновесием полной системы только при принадлежности соответствующей области.

В физическом варианте $z=C$. В смешанном варианте $z=\Phi(C)$ и $C=\Phi^{-1}(z)$; скрытая переменная $q$ не отождествляется с рецепторной переменной $h$. В используемой смешанной семье при $u=C/(0.5\,\text{мкМ})$:
\begin{equation}
\begin{gathered}
\Phi(C)=w_1\frac{u^n-1}{n}+w_2\frac{e^{k(u-1)}-1}{k}+w_3\ln u,\\
(w_1,w_2,w_3)=\frac{(e^p,e^e,1)}{e^p+e^e+1}.
\end{gathered}
\end{equation}
При $k=0$ второе слагаемое понимается по непрерывности. Для сохранённых положительных $n,k$ производная положительна при $C>0$, и обратное преобразование однозначно на принятом физическом интервале. PWA-свойство относится к $(z,q)$: поле в физических координатах после нелинейного преобразования обычно не является кусочно-аффинным.

\subsection{Локальные переходы HLB1 и HLB2}
При пересечении равновесием границы меняется региональный якобиан. Возникновение цикла в непрерывной кусочно-гладкой системе в этом случае может иметь Hopf-подобный характер~\cite{Simpson2022,SimpsonMeiss2007}. Рассматриваются два варианта: фокус--фокус (HLB1) и узел--фокус (HLB2). Они отличаются от гладкой бифуркации Хопфа и требуют проверки условий пересечения и локального отображения возврата.

На рис.~\ref{fig:ru-hlb} показаны конкретные безразмерные примеры $\dot z=y-F(z)$, $\dot y=\delta-z$ с порогом $z=0$. Для HLB1 наклоны $F$ равны $(1,-0.45)$, для HLB2 --- $(3,-0.45)$. При $\delta=0.12$ независимое интегрирование даёт устойчивые циклы с периодами соответственно $6.6401$ и $7.1212$, а логарифмы нетривиальных множителей Флоке равны $-0.9970$ и $-5.5471$. Это иллюстративные системы, а не параметры моделей, подобранных к Li--Rinzel.
\rufigure{figure3_hlb}{Примеры локальных циклов HLB1 и HLB2. Показаны нульклины, направления поля и устойчивые периодические орбиты. Подписи областей характеризуют соответствующие аффинные подсистемы; действительное равновесие обозначено точкой. Координаты и периоды безразмерны.}{fig:ru-hlb}{}

\subsection{Большие циклы и область применимости теорем}
Большими называются циклы, пересекающие обе границы трёхзонной системы. На рис.~\ref{fig:ru-bc} показаны два способа организации такой орбиты: BC1 с фокусной левой областью и BC2 с узловой левой областью. В обоих примерах средняя область фокусная неустойчивая, правая --- фокусная устойчивая. Наклоны $F$ равны $(1,-0.45,1)$ и $(3,-0.45,1)$ соответственно, пороги $\pm1$. Периоды равны $6.4218$ и $6.5934$, логарифмы множителей --- $-2.8627$ и $-5.8549$. Обозначения BC1 и BC2 относятся к двум конфигурациям, а не к двум одновременно существующим циклам одной системы.

Llibre, Ponce и Valls~\cite{LlibrePonceValls2015} доказывают существование двух циклов при специальных предпосылках для непрерывных трёхзонных систем. Chen и соавторы~\cite{ChenFengJiaMeng2026} устанавливают максимальное число два для системы Льенара с глобально линейной восстанавливающей частью и непрерывной кусочно-линейной функцией с двумя изломами. Оценка сверху не доказывает существование двух циклов и сама по себе не классифицирует BC1/BC2.

Для применения последнего результата к~\eqref{eq:pwa-ru} требуется допустимое преобразование к изученному классу. Исключение $q$ внутри области даёт
\begin{equation}
\begin{aligned}
\ddot z+(a n_j+b)\dot z&+ab[N(z,I)-H(z,I)]\\
 &=-a\alpha_I\dot I.
\end{aligned}
\end{equation}
При постоянном $I$ восстанавливающая часть в общем случае лишь кусочно-аффинна, с различными наклонами $ab(n_j-h_j)=D_j$. Поэтому без отдельного доказательства приведения теорема 2026 года не даёт общей границы в два цикла для всех наших подборов. Для фактически найденных циклов далее используется непосредственная проверка отображений возврата.
\rufigure{figure4_large_cycles}{Большие устойчивые циклы BC1 и BC2, пересекающие обе границы. В BC1 левая аффинная подсистема фокусная, в BC2 --- узловая. Это два безразмерных примера с разными параметрами; рисунок не показывает сосуществование двух циклов при одном наборе коэффициентов.}{fig:ru-bc}{}

\subsection{Аналитическое распространение состояния}
Внутри зоны обозначим $X=(z,q)^\mathsf{T}$. Уравнения имеют вид $\dot X=A_jX+B I+c_j$. При постоянном входе
\begin{equation}
\begin{aligned}
X(t+s)&=e^{A_js}X(t)\\
 &\quad+\int_0^s e^{A_j(s-u)}(BI+c_j)\,du.
\end{aligned}
\end{equation}
Если между входными событиями $\dot I=-(I-I_b)/\tau$, добавим к состоянию $I$ и единицу:
\begin{equation}
\begin{gathered}
Y=\begin{pmatrix}X\\I\\1\end{pmatrix},\qquad
\dot Y=\mathcal A_jY,\\
\mathcal A_j=\begin{pmatrix}A_j&B&c_j\\0&-1/\tau&I_b/\tau\\0&0&0\end{pmatrix},\\
Y(t+s)=e^{\mathcal A_js}Y(t).
\end{gathered}
\end{equation}
Формула не требует обратимости $A_j$ и применима до первого выхода из области. Время пересечения порога находится решением скалярного уравнения для $z(t+s)$; после переключения используется матрица следующей области. Входной скачок IP$_3$ обрабатывается отдельно. При постоянном входе цикл получается согласованием региональных переходов, а его устойчивость --- производной отображения возврата. Для непрерывной плоской системы при трансверсальных пересечениях $\log M=\sum_jT_jt_j$. Условие $M=1$ характеризует нейтральную орбиту; для складки дополнительно проверяются невырожденность и зависимость возврата от входа.

\section{Подбор и сравнение с Li--Rinzel}
\subsection{Протокол сравнения}
Для каждого из трёх режимов сравниваются физические и смешанные координаты, с фиксацией верхней складки циклов и без неё. Использованы семь обучающих и четыре проверочных значения входа. Сравниваются абсолютные ошибки минимума и максимума $C$, периода и RMSE формы цикла. При вычислении ошибки формы циклы нормированы по фазе и совмещены по максимуму; эта метрика не измеряет накопление фазовой ошибки при длительном внешнем входе. Все двенадцать основных проверочных строк содержат четыре найденных цикла из четырёх.

Фиксация складки и фиксация входа пересечения равновесием порога --- различные ограничения. В вариантах со складкой согласовывается нейтральная орбита, пересекающая оба порога, при входе верхней складки Li--Rinzel. Для AM соответствующая эталонная орбита не пересекает нижний порог. Таким образом, ограничение на маршрут PWA-орбиты не воспроизводит маршрут эталона и может приводить к компромиссу между глобальным условием и локальной точностью.
\todo{Перед подачей свести точные веса целевой функции, диапазоны параметров, бюджеты оптимизации и критерии остановки из исходных запусков. Исторические варианты используют различающиеся ограничения скрытой координаты; их нельзя представлять как полностью контролируемую абляцию одного условия.}

\subsection{Точность периодов и формы циклов}
Таблица~\ref{tab:ru-errors} составлена непосредственно по проверочной части сохранённых данных. Освобождение складки улучшает физический AM: ошибка периода уменьшается с $3.167$ до $0.205$ с, RMSE формы --- с $0.06347$ до $0.01764$ мкМ. В AFM соответствующие ошибки изменяются с $1.803$ до $0.112$ с и с $0.12832$ до $0.03532$ мкМ. Для физического FM освобождение складки уменьшает ошибку максимума с $0.21791$ до $0.12754$ мкМ, но не улучшает период и форму. Следовательно, эффект ограничения зависит от режима и выбранной метрики.

При свободной складке смешанные координаты дают меньшую RMSE формы, чем физические: $0.01115$ против $0.01764$ мкМ в AM, $0.01054$ против $0.03532$ мкМ в AFM и $0.06872$ против $0.20460$ мкМ в FM. Однако период AFM в смешанном варианте несколько менее точен. Эти результаты показывают полезность нелинейного считывания кальция, но не универсальное улучшение каждой характеристики.
\begin{table}[p]\centering\small
\caption{Ошибки на четырёх проверочных входах. Для каждой строки найдены 4/4 цикла. Ошибки экстремумов и формы указаны в мкМ, периода --- в секундах.}\label{tab:ru-errors}
\resizebox{\linewidth}{!}{\begin{tabular}{llcrrrr}\toprule
Режим & Координаты & Складка & $E_{\min}$ & $E_{\max}$ & $E_T$ & RMSE\\\midrule
AM & Физич. & фикс. & 0.0292 & 0.0740 & 3.1672 & 0.0635\\
AM & Физич. & своб. & 0.0220 & 0.0144 & 0.2054 & 0.0176\\
AM & Смеш. & своб. & 0.0124 & 0.0104 & 0.1076 & 0.0112\\
AM & Смеш. & фикс. & 0.0302 & 0.0731 & 3.0236 & 0.0632\\
AFM & Физич. & фикс. & 0.0039 & 0.1605 & 1.8028 & 0.1283\\
AFM & Физич. & своб. & 0.0163 & 0.0547 & 0.1123 & 0.0353\\
AFM & Смеш. & своб. & 0.0112 & 0.0032 & 0.1686 & 0.0105\\
AFM & Смеш. & фикс. & 0.0066 & 0.0373 & 0.9280 & 0.0695\\
FM & Физич. & фикс. & 0.0024 & 0.2179 & 0.6534 & 0.2009\\
FM & Физич. & своб. & 0.0049 & 0.1275 & 0.7440 & 0.2046\\
FM & Смеш. & своб. & 0.0025 & 0.0617 & 0.3828 & 0.0687\\
FM & Смеш. & фикс. & 0.0051 & 0.2759 & 0.4085 & 0.1713\\
\bottomrule\end{tabular}}\end{table}


Сравнение ветвей и периодов приведено на рис.~\ref{fig:ru-periods} и~\ref{fig:ru-branches}. Складки без ограничения являются результатом последующего поиска, а не навязанными целями. Отсутствие ветви в конечном поиске не доказывает её отсутствия в системе. Для свободного FM в этих рисунках используется более ранний архивный кандидат: окончательный замороженный файл не был доступен при построении. Поэтому графические ветви и агрегированные ошибки финального FM нельзя считать описанием одного и того же файла модели.
\rufigure{figure5_period_comparison}{Периоды Li--Rinzel и смешанных PWA-моделей: эталон, подбор с фиксированной верхней складкой и подбор без её фиксации. Заполненные и пустые маркеры --- устойчивые и неустойчивые циклы, ромбы --- складки. Поиск ветвей конечен. В свободном FM показан ранний архивный кандидат, а не окончательная модель из таблицы ошибок.}{fig:ru-periods}{}
\rufigure{figure6_bifurcation_comparison}{Равновесия и экстремумы циклов в тех же вариантах. Сплошная линия --- устойчивые равновесия, пунктир --- неустойчивые, включая седла. Ограничение фиксированной складки AM требует пересечения обоих порогов, что отличается от маршрута эталонной орбиты. Происхождение свободного FM совпадает с рис.~\ref{fig:ru-periods}.}{fig:ru-branches}{}

\section{Вычислительная эффективность и проверка в NEST}
Модели включаются в нейроастроцитарную сеть NEST~\cite{GewaltigDiesmann2007,Jiang2025}. Физический кальций восстанавливается из внутреннего состояния и используется для вычисления SIC. Поэтому стоимость сетевой симуляции включает распространение состояния, локализацию границ, обратное преобразование, SIC, нейронную динамику и регистрацию данных. Аналитический региональный поток не меняет автоматически шаговый планировщик NEST.

В первом тесте сеть содержит 10 нейронов и 90 астроцитов, сопоставленных направленным межнейронным связям полносвязной сети. Моделируются 60 с при шаге $0.1$ мс и 16 потоках. Для каждого варианта выполнены три повтора с одинаковым seed; анализируется медиана времени \texttt{nest.Run}. При интервале регистрации 10 мс ускорение физических моделей составляет $2.78$, $3.16$ и $2.65$ для AM, AFM и FM; смешанных --- $2.94$, $2.39$ и $2.71$. По двум интервалам регистрации, 10 и 1000 мс, общий диапазон составляет $2.39$--$3.16$.

В отдельной FM-сети со 100 нейронами число астроцитов увеличивалось от 100 до 1600. Для физической модели отношение времён возрастает с $1.58$ до $5.19$, для смешанной --- с $1.47$ до $3.42$. При 1600 астроцитах соответствующие медианы LR/PWA равны $24.86/4.79$ с и $24.78/7.25$ с. Для этого опыта доступны медианы, но не сырые три повтора. Динамика также различается: для первой записанной траектории в таблице указаны три пика LR против двух PWA. Ускорение не следует интерпретировать как равноточную замену астроцита.
\rufigure{figure7_nest_speed}{Вычислительная эффективность в NEST. Слева --- отношение медианных времён в сети с 10 нейронами и 90 астроцитами при разных интервалах регистрации. В центре и справа --- ускорение и абсолютные времена в отдельной FM-сети со 100 нейронами при изменении числа астроцитов. SIC включён. Времена отражают полную сетевую динамику и не выделяют стоимость преобразования координат.}{fig:ru-speed}{}

Физические и смешанные варианты подбирались независимо. Их различие по времени включает изменение траекторий, числа переключений и сетевой активности. Следовательно, имеющиеся данные оценивают совместный компромисс точности и скорости, но не причинный эффект одного обратного преобразования. Для выделения его стоимости нужен отдельный тест считывания одной и той же последовательности состояний с SIC и без SIC при одинаковом расписании обращений.
\todo{Дополнить окончательный протокол точной версией NEST, CPU, параметрами SIC, графом связей, начальными состояниями и идентификаторами финальных файлов. Сохранённые SHA256 кандидатов и исходные повторные времена доступны в таблицах проекта; измерения прежних многозонных моделей не переносятся на текущие трёхзонные.}

\section{Обсуждение}
Три области позволяют организовать устойчивое состояние, промежуточную неустойчивую фокусную динамику и восстановление устойчивости при больших входах. Свобода типа левой области даёт две локальные конфигурации и два варианта больших циклов. При этом тип аффинной подсистемы, действительное равновесие и существование цикла являются различными свойствами и проверяются отдельно.

Близость периодов и амплитуд не равносильна сохранению всех бифуркаций. Особенно чувствительны нижний FM-переход, тонкие локальные циклы и маршрут верхней складки AM. Высокоточные локальные кандидаты AFM/FM в сохранённых материалах не имеют полной независимой проверки и разрешения по параметру; они не используются здесь как доказательство воспроизведения наблюдаемой ветви Li--Rinzel. Общая граница числа циклов также требует проверки класса системы, а не только количества зон.

Нелинейное восстановление кальция расширяет выразительность компактной модели и улучшает форму ряда ответов. Его цена становится существенной при частом считывании состояния, поэтому сравнение только изолированных решателей было бы недостаточным. Сетевые результаты демонстрируют практический выигрыш в рассмотренных конфигурациях, сохраняя видимым расхождение динамики. Для более широкого применения требуется проверка на новых временных входах и сетевых режимах без повторного подбора.

\section{Заключение}
Непрерывный трёхзонный осциллятор представляет компактное феноменологическое описание выбранных кальциевых режимов Li--Rinzel. Региональная аффинность обеспечивает аналитическое распространение при постоянном и экспоненциально релаксирующем IP$_3$. Сравнение сохранённых подборов выявляет зависимый от режима компромисс между точностью траекторий и фиксацией верхней складки. Смешанные координаты уменьшают ошибку формы циклов, а сетевые тесты NEST показывают сокращение времени расчёта при сохраняющейся ошибке модели. Полная бифуркационная эквивалентность, общая оценка числа циклов для всех подборов и отдельная стоимость преобразования координат остаются самостоятельными вопросами.

\section*{Доступность данных и кода}
Код, параметры, рисунки и исходные таблицы размещены в репозитории проекта. Для каждого сравнения требуется сохранять принадлежность конкретному кандидату и протоколу; окончательную архивную ссылку следует добавить перед подачей.
\section*{Вклад авторов, финансирование и конфликт интересов}
Сведения заполняются авторами перед подачей рукописи.
